% ============================================================================
% CAPÍTULO 2 — MARCO TEÓRICO
% Objetivo: derivar la ecuación de movimiento de la cáscara desde el formalismo.
% Resultado que DEBE contener:
%   · Condiciones de empalme de Israel (salto de la curvatura extrínseca).
%   · La ecuación maestra GENERAL en términos de f_±(R) — la fuente de verdad:
%       sqrt(Rdot^2 + f_-) - sqrt(Rdot^2 + f_+) = m/R
%   · La reducción a ODE para R(tau) via Birkhoff (por qué NO hace falta NR).
% Nota de honestidad para vos: acá va la derivación con los índices cuidados
%   (intrínseco vs. ambiente) — es donde históricamente aparecían los errores.
% Refs: Israel 1966; Poisson §3.7-3.8.
% ============================================================================
\chapter{Marco teórico: condiciones de empalme de Israel}
\label{cap:marco}

\section{Cáscaras delgadas y curvatura extrínseca}
% TODO: hipersuperficie tipo tiempo, n^mu espacelike, salto de K_ij.

\section{Ecuación de movimiento general}
\label{sec:eom-general}
% La condición de empalme para una cáscara de masa en reposo m(R):
\begin{equation}
    \sqrt{\Rdot^2 + f_-(R)} \;-\; \sqrt{\Rdot^2 + f_+(R)} \;=\; \frac{m(R)}{R}.
    \label{eq:maestra}
\end{equation}
% TODO: despejar Rdot^2 en general (resta de cuadrados) y comentar que este es
%   el núcleo que consumen integrador, benchmark y mapa.

\section{Reducción a una ecuación diferencial ordinaria}
% TODO: teorema de Birkhoff fija la geometría a ambos lados -> problema = ODE
%   para R(tau). Argumento explícito de por qué no se requiere NR 3+1.
